\scanpage{92}
\mspage{3.}

de $\pi_1^t$ par $I'_0$, soit \deleted{\unsure{$\Pi'_2$, un}} $\Pi'$, et nous \deleted{\unsure{dev}} \unsure{aurons|avons|donc aussi} un \uline{épimorphisme}
\[
 \Pi'\longrightarrow\pi_1^t(\unsure{\mathfrak X|\mathscr X|\mathcal X}/D)
\]
\unsure{prouvons} que c'est un \uline{isomorphisme}. ]

Ceci revient à vérifier ceci : étant donné un hom surjectif
\[
 \mu_0\longrightarrow J_0\longrightarrow0
\]
qui s'annule sur l'image de $\pi_2^t(U)$, \deleted{\unsure{il existe}} $J_0$ apparaît comme groupe d'inertie dans un \unsure{revêtem|revêt.} \deleted{\unsure{tame}} modéré de $\mathfrak X/X_0$.

L'hypothèse signifie aussi que l'image de $\beta\in H^2(\unsure{U|X|\mathcal U},\mu)$ dans $H^2(\widetilde U^t,J_0)$ est nulle ; elle devient donc nulle déjà \deleted{\unsure{ap}} à un niveau fini [\deleted{NB.} \uline{par définition}, on définira la coh de $\widetilde U^t$ comme une \uline{$\lim$} .... si $\unsure{U|\mathfrak X}$ n'est pas noeth \unsure{et} mais seulement loc. noeth .... ]. Donc quitte à passer à un \unsure{revêtem} modéré fini de \deleted{\unsure{variété} \unsure{tame}, on peut} \inserted{\deleted{$\mathfrak X$ \unsure{relativ}} /($D_i$ \unsure{laissés|lisses},} \deleted{on peut} \deleted{\unsure{laisse}}
\margin{\unsure{Composé} de rev. tames est tame}
% archive sheet number: 91
